\begin{ledgroupsized}[r]{120mm}
\footnotesize
\pstart
\noindent\textbf{Überlieferung:}
\pend
\end{ledgroupsized}
\begin{ledgroupsized}[r]{114mm}
\footnotesize
\pstart \parindent -6mm
\makebox[6mm][l]{\textit{L}}%
Konzept:
LH~XXXV~13, 3~Bl.~242. 
Ein Blatt~2\textsuperscript{o};
Wasserzeichen		%
in der Blattmitte;
linker Rand beschnitten.
Eineinviertel Seiten.
\pend 
\end{ledgroupsized}
\vspace{5mm}
\begin{ledgroup}
\footnotesize
\pstart
\noindent%
\textbf{Datierungsgründe:}
Die Datierung des vorliegenden Konzeptes stützt sich auf eine Anspielung auf den \textit{vis viva}-Streit mit 
\indexN{\textls{Papin} (Papinus), Denis 1647\textendash1713}Denis Papin.
Leibniz beweist in N.\;\ref{58300} die Erhaltung der verschiedenen Formen der \textit{vis} im Fall eines schiefen Stoßes zweier gleicher Körper auf einen dritten ruhenden. 
Da die stoßenden Körper dabei zum Stillstand kämen, werde die gesamte Kraft an den dritten übertragen.
Leibniz merkt an, dass dieser Beweis auch gegen \indexN{\textls{Papin} (Papinus), Denis 1647\textendash1713}Papin einzusetzen sei, weil daraus die 
von diesem geleugnete
Übertragbarkeit der gesamten \textit{vis motrix} eines Körpers an einen ruhenden Körper geringerer Masse folge (\refpassage{35_13_03_242_7a}{35_13_03_242_7b}).
Leibniz hatte in seiner ersten Streitschrift gegen \indexN{\textls{Papin} (Papinus), Denis 1647\textendash1713}Papin
(\cite{02040}\textit{De causa gravitatis}, \cite{01023}\textit{AE} Mai 1690, S.\;228\textendash239)
die uneingeschränkte Übertragbarkeit der \textit{vis motrix} beim elastischen Stoß vorausgesetzt.
Im Aufsatz 
\cite{00537}\textit{Mechanicorum de viribus motricibus sententia} (\cite{01023}\textit{AE} Januar 1691, S.\;6\textendash13, bes.\ S.\;9\textendash11)
bestritt \indexN{\textls{Papin} (Papinus), Denis 1647\textendash1713}Papin 
mit Verweis auf die
\indexN{\textls{Huygens} (Hugenius, Ugenius, Hugens, Huguens), Christiaan 1629\textendash1695}Huygens'schen Regeln des geraden Stoßes
diese Möglichkeit,
besonders im genannten Fall des Stoßes eines größeren Körpers auf einen kleineren ruhenden.
\indexN{\textls{Papin} (Papinus), Denis 1647\textendash1713}Papins Aufsatz, dem die Anspielung auf \refpassage{35_13_03_242_7a}{35_13_03_242_7b} gilt, bildet daher einen Terminus post quem für die Entstehung von N.\;\ref{58300}.
Einen Terminus ante quem bietet wiederum Leibnizens
\cite{02146}Brief an \indexN{\textls{Bernoulli}, Johann 1667\textendash1748}Johann Bernoulli vom 28.~Januar (7.~Februar) 1696,
in dem er 
(wie auch in N.\;\ref{58300})
erwägt, sein neues Gedankenexperiment künftig gegen \indexN{\textls{Papin} (Papinus), Denis 1647\textendash1713}Papin einzusetzen
(\textit{LSB} III,~6 N.\;202, hier S.\;648.12\textendash16).
Erst in einem \cite{02073}Schreiben vom 16. (26.) Januar 1698 teilt Leibniz sein Argument \indexN{\textls{Papin} (Papinus), Denis 1647\textendash1713}Papin selbst mit 
(\textit{LSB} III, 7, N.\;177, hier S.\;724).
\pend
\pstart
Das für den Zeitraum Oktober 1690 bis 
1696 mehrfach belegte Wasserzeichen bestätigt die angegebene Datierung. %
\pend
\end{ledgroup}
\selectlanguage{latin}
\vspace{8mm}
\pstart%
\normalsize%
\noindent%
\href{\Scan{LH_35_13_03_242r}}{\lbrack242~r\textsuperscript{o}\rbrack}
\pend
\nobreak\begin{edfigure}[0pt][1em]    	%
\includegraphics[scale=1]{%
gesamttex/edit_VIII-4/images/LH_35_13_03_242_d1_242r.pdf}
\edcaption[0pt]{\edtext{\lbrack\textit{Fig.~1}\rbrack}{%
\lemma{\lbrack\textit{Fig.~1}\rbrack}%
\Cfootnote[nonum]{Ein gestrichener Entwurf zum Diagramm wird nicht wiedergegeben.}}}
\end{edfigure}
\clearpage
\pstart
\noindent
Sint globi aequales\indexS{globi aequales} \textit{A}, \textit{B}, \textit{C}, \textit{D}, incurratque globus \textit{A} %
\edtext{(solus motus)}{\lemma{}\Bfootnote{(solus motus) \textit{erg.~L}}}
in globum \textit{B}, tangentem globos \textit{C} et \textit{D}. Videndum an %
\edtext{post ictum\indexS{ictus}}{\lemma{}\Bfootnote{post ictum \textit{erg.~L}}}
quiescant globi \textit{A} et \textit{B}, solumque moveantur \textit{C} et \textit{D}. %
\edtext{Hoc posito}{\lemma{}\Bfootnote{Hoc posito \textit{erg.~L}}} 
\edtext{celeritas ipsius \textit{A} sit 1 et celeritas ipsorum \textit{C} et \textit{D} erit %
\edlabel{35_13_03_242_1a}%
\edtext{}{{%
\xxref{35_13_03_242_1a}{35_13_03_242_1b}}%
\lemma{$1:\sqrt{2}$.}%
\Bfootnote{%
\textit{(1)}~Celeritas %
\textit{(2)}~Erit \textit{L} %
\textit{(3)}~Celeritas centri gr. %
\textit{(4)}~\bars{Nam 1 \textit{streicht Hrsg.}} %
\textit{(5)}~Ponamus~\textit{L}}}%
$1:\sqrt{2}$.}%
{\lemma{}\Afootnote{\textit{Oberhalb des Textes}: $\displaystyle\frac{2}{xx}=1$. \;$x=\lbrack\sqrt2\rbrack$\textsuperscript{\lbrack a\rbrack}. Ergo $1\cdot 1A=\displaystyle\frac{1}{\sqrt2}\cdot \displaystyle\frac{1}{\sqrt2}B+\displaystyle\frac{1}{\sqrt2}\cdot \displaystyle\frac{1}{\sqrt2}\cdot C$\lbrack,\rbrack\ si $A=B=C$
\\[1.5ex] %
{\footnotesize \textsuperscript{\lbrack a\rbrack} $1:\sqrt2$ \textit{L ändert Hrsg.}}
}}%
\pend
\begin{edfigure}[.8em][.5em]    	 %
\includegraphics[scale=1]{%
gesamttex/edit_VIII-4/images/LH_35_13_03_242_d2_242r.pdf} 
\edcaption[0.25em]{\lbrack\textit{Fig.~2, gestr.}\rbrack}
\end{edfigure}
\pstart
Ponamus\edlabel{35_13_03_242_1b}
\textit{\ind{1}B\ind{1}C\ind{1}D} esse triangulum rectangulum %
\edtext{isosceles. Et celeritas ipsius \textit{M} centri gravitatis\indexS{celeritas centri gravitatis}\indexS{centrum gravitatis}}{\lemma{isosceles}\Bfootnote{%
\textit{(1)}~, erunt %
\textit{(2)}~. Et %
\textit{(a)}~centrum gravitatis %
\textit{(b)}~\bars{\textit{M} \textit{streicht Hrsg.}} centrum gravitatis %
\textit{(c)}~celeritas ipsius \textit{M} centri gravitatis~\textit{L}}} 
ipsorum \textit{C} et \textit{D} post ictum
\edtext{erit %
\rule[-1.8ex]{0pt}{5.4ex}%
$\displaystyle\frac{1}{2}$. Porro}{\lemma{erit $\displaystyle\frac{1}{2}$.}\Bfootnote{%
\textit{(1)}~Jam %
\textit{(2)}~Si duo sint %
\textit{(a)}~\textit{L} et \textit{M}, %
\textit{(b)}~puncta \textit{L} et \textit{M}, sitque \textit{M} ut 3, et \textit{L} ut 1 aequalia et quiescente \textit{L} moveatur \textit{M}, %
\textit{(aa)}~centrum %
\textit{(bb)}~centri gravitatis progressus est ut %
\textit{(aaa)}~1 %
\textit{(bbb)}~\rule[-2.1ex]{0pt}{5.4ex}%
$\displaystyle\frac{1}{2}$
\textit{(3)}~Generaliter duo movean %
\textit{(4)}~Porro~\textit{L}}}
generaliter si duo moveantur in eadem recta \textit{L} et \textit{M}, sitque ipsius \textit{L} velocitas \textit{c} et ipsius \textit{M} velocitas %
\edtext{\lbrack\textit{e}\rbrack,}{\lemma{}\Bfootnote{\textit{E} \textit{L ändert Hrsg.}}}
erit centri gr.\ velocitas\indexS{velocitas centri gravitatis} $Lc+eM\lrcorner:\llcorner l+m$\,, hinc si unum \textit{L} quiescat, erit centri
\edtext{gravita\lbrack ti\rbrack s}{%
\lemma{}%
\Bfootnote{%
gravitas %
\textit{L ändert Hrsg.}}}
velocitas\indexS{velocitas centri gravitatis} $eM:\llcorner l+m$. Ergo si \textit{M} sit centrum gravitatis\indexS{centrum gravitatis} ipsorum \textit{C} %
\edtext{et \textit{D}; at}{\lemma{et \textit{D};}\Bfootnote{%
\textit{(1)}~\bars{et \textit{streicht Hrsg.}} \textit{M} ipsorum %
\textit{(2)}~at~\textit{L}}} 
\textit{A} %
\textit{L} sit centrum grav.\ ipsorum
\edtext{et \textit{B}, utique}{\lemma{et \textit{B},}\Bfootnote{%
\textit{(1)}~et, %
\textit{(2)}~utique~\textit{L}}} 
\textit{M} erit ut 2, et \textit{L} etiam, %
\edtext{ob duos}{\lemma{}\Bfootnote{ob \bars{ob \textit{streicht Hrsg.}} duos~\textit{L}}} 
utrobique globos, et \textit{e},
\edtext{velocitas ipsius}{\lemma{velocitas}\Bfootnote{%
\textit{(1)}~centri grav. %
\textit{(2)}~ipsius~\textit{L}}}
\textit{M} est 
\rule[-1.8ex]{0pt}{5.4ex}%
$\displaystyle\frac{1}{2}$.
Ergo centri gravitatis communis\indexS{centrum gravitatis commune} 4 globorum velocitas\indexS{velocitas centri gravitatis} post
\edtext{ictum $eM:\llcorner l+m$}{%
\lemma{ictum}%
\Bfootnote{%
\textit{(1)}~erit %
\textit{(2)}~$eM:\llcorner l+m$~\textit{L}}}
erit%
\linebreak
\edlabel{35_13_03_242_2a}%
\edtext{}%
{{\xxref{35_13_03_242_2a}{35_13_03_242_2b}}%
\lemma{%
$\displaystyle\frac{1}{2} \cdot 2:\llcorner2+2=\displaystyle\frac{1}{4}$.}%
\Bfootnote{\rule[-1.8ex]{0pt}{5.4ex}%
\textit{(1)}~Sed ante ictum %
\textit{(2)}~At quaeramus ejusdem %
\textit{(3)}~Quaeramus~\textit{L}}}%
\rule[-1.4ex]{0pt}{5.ex}%
$\displaystyle\frac{1}{2} \cdot 2:\llcorner2+2=\displaystyle\frac{1}{4}$.
\pend
\clearpage
\pstart %
Quaeramus%
\edlabel{35_13_03_242_2b}
jam et quae fuerit ante 
\edtext{ictum.\indexS{ictus} \textit{A}}{%
\lemma{ictum.}%
\Bfootnote{%
\textit{(1)}~Tunc centrum %
\textit{(a)}~velocitatis %
\textit{(b)}~gravitatis %
\textit{(aa)}~ipsi %
\textit{(bb)}~ipsius \textit{A} erit \textit{A} seu 1 %
\textit{(aaa)}~, ipsorum trium
\textit{(bbb)}~at trium quod exprimatur per \textit{N}
\textit{(2)}~\textit{A}~\textit{L}}}
sit centrum gravitatis\indexS{centrum gravitatis} 
\edtext{ipsius \textit{A}, quod}{%
\lemma{ipsius \textit{A},}%
\Bfootnote{%
\textit{(1)}~et \textit{N} sit centrum gr.\ trium %
\textit{(2)}~quod~\textit{L}}}
movetur, et \textit{N} sit centrum gravitatis trium\indexS{centrum gravitatis trium} \textit{B}, \textit{C},
\edtext{\textit{D}, quod quiescit. Ergo}{%
\lemma{\textit{D},}%
\Bfootnote{%
\textit{(1)}~qu %
\textit{(2)}~jam %
\textit{(3)}~, quod quiescit. Ergo~\textit{L}}}
velocitas ipsius \textit{N} est 0, et ipsius \textit{A} est 1, fiet velocitas centri gravitatis\indexS{velocitas centri gravitatis} omnium communis ante ictum $1A+0N\lrcorner:\llcorner A+N$. Est autem \textit{A} ut 1, et \textit{N} ut 3, et fiet \rule[-2ex]{0pt}{5.4ex}%
$\displaystyle\frac{1}{4}$ ut ante. Haec ergo %
\edtext{recte. Servatur}{\lemma{recte.}\Bfootnote{%
\textit{(1)}~Et vis serv %
\textit{(2)}~Servatur~\textit{L}}} 
enim eadem vis communis\indexS{vis communis}
\edtext{quae ante concursum.\indexS{concursus} Vis totalis\indexS{vis totalis} $1\cdot 1\cdot A=1$. 
Vis communis\indexS{vis communis} 
\rule[-2.4ex]{0pt}{6ex}%
$\displaystyle\frac{A+B+C+D}{16}=\displaystyle\frac{1}{4}$. 
Propria\indexS{vis propria} \rule[-1.8ex]{0pt}{5.4ex}$\displaystyle\frac{3}{4}$. Videamus}{%
\lemma{}%
\Bfootnote{communis
\bars{%
\textit{(1)}~quae ante concursum 
\textit{(2)}~quae ante concursum 
\textit{(a)}~\rule[-1.8ex]{0pt}{5.4ex}%
$\displaystyle\frac{A+B+C+D}{16}$ $=\displaystyle\frac{1}{4}$ 
\textit{(b)}~. Vis totalis $1\cdot 1\cdot A=1$. Vis communis 
\rule[-1.8ex]{0pt}{5.4ex}%
$\displaystyle\frac{A+B+C+D}{16}=\displaystyle\frac{1}{4}$. Propria $\displaystyle\frac{3}{4}$ 
\textit{erg.}} 
.~Videamus~\textit{L}}} 
nunc an etiam eadem vis propria\indexS{vis propria} servetur.
\pend
\pstart 
Vis propria\indexS{vis propria} ante %
\edtext{concursum sic}{\lemma{concursum}\Bfootnote{%
\textit{(1)}~erat %
\textit{(a)}~$\displaystyle\frac{1}{6}\cdot$ 
\textit{(b)}~$\displaystyle\frac{1}{9}A +$ %
\textit{(2)}~sic~\textit{L}}} 
habebitur, %
\edtext{si respectiva celeritas\indexS{celeritas respectiva} secta}{\lemma{si}\Bfootnote{%
\textit{(1)}~appropinquandi %
\textit{(2)}~respectiva celeritas %
\textit{(a)}~nempe 1, %
\textit{(aa)}~dividatur in duas partes %
\textit{(bb)}~secetur in ratione reci %
\textit{(b)}~secta~\textit{L}}} %
in ratione \textit{A} ad 
\edtext{\textit{N}, inter}{\lemma{\textit{N},}\Bfootnote{%
\textit{(1)}~reciproce %
\textit{(2)}~inter~\textit{L}}} 
ipsa reciproce 
\edtext{distribuatur. Erat}{%
\lemma{distribuatur.}%
\Bfootnote{%
\textit{(1)}~Et %
\textit{(2)}~Erat~\textit{L}}}
autem celeritas appropinquandi\indexS{celeritas appropinquandi} 1. Secanda in 
\rule[-1.8ex]{0pt}{6ex}%
$\displaystyle\frac{1}{4}$ et $\displaystyle\frac{3}{4}$ 
et ipsi \textit{A} ascribetur celeritas 
$\displaystyle\frac{3}{4}$, 
ipsi \textit{N} celeritas 
\rule[-1.8ex]{0pt}{6ex}%
$\displaystyle\frac{1}{4}$. 
Et vis respectiva %
\edtext{ante concursum}{\lemma{}\Bfootnote{ante concursum \textit{erg.~L}}} 
erit $9A+1N\lrcorner:16=12:16=\displaystyle\frac{3}{4}$.%
\rule[-1.8ex]{0pt}{5.4ex}
Videamus jam quae sit post concursum.\indexS{concursus} Est %
\edtext{autem bipertita, scilicet vis respectiva\indexS{vis respectiva}}{\lemma{autem}\Bfootnote{%
\textit{(1)}~duplex, vis respectiva %
\textit{(2)}~bipertita, scilicet vis respectiva~\textit{L}}} 
inter \textit{L} et \textit{M}, et vis respectiva\indexS{vis respectiva} inter \textit{C} et \textit{D}. Et quidem celeritas respectiva\indexS{celeritas respectiva} inter \textit{L} et \textit{M}, est celeritas 
\edtext{ipsius \textit{M},}{%
\lemma{ipsius}%
\Bfootnote{%
\textit{(1)}~\textit{N}, nempe %
\textit{(2)}~\textit{M},~\textit{L}}}
nempe \rule[-1.8ex]{0pt}{5.4ex}%
$\displaystyle\frac{1}{2}$. Jam \textit{L} et \textit{M} sunt 
\edtext{aequalia pondere,}{%
\lemma{aequalia}%
\Bfootnote{%
\textit{(1)}~pondera, %
\textit{(2)}~pondere,~\textit{L}}}
quia ambo centra grav.\ 2 globorum. Ergo celeritas respectiva distributa\indexS{celeritas respectiva distributa} est 
\rule[-1.8ex]{0pt}{5.4ex}%
$\displaystyle\frac{1}{4}$. 
Et vis respectiva\indexS{vis respectiva} inter \textit{L} et \textit{M} erit 
\rule[-1.8ex]{0pt}{5.4ex}%
$\displaystyle\frac{2}{16}+\displaystyle\frac{2}{16}=\displaystyle\frac{1}{4}$. 
Restat vis respectiva\indexS{vis respectiva} inter %
\edtext{\textit{C} et \textit{D}. Celeritas}{\lemma{\textit{C} et \textit{D}}\Bfootnote{%
\textit{(1)}~. Est %
\textit{(2)}~quae est %
\textit{(3)}~. Et celeritas cujusque respectiva est 1. Ergo vis %
\textit{(a)}~eorum respectiva est %
\textit{(b)}~cujusque est \rule[-1.8ex]{0pt}{5.4ex}%
$\displaystyle\frac{1}{4}$ %
\textit{(c)}~eorum distributa est $\displaystyle\frac{1}{2}$ %
\textit{(4)}~. Celeritas~\textit{L}}} 
eorum respectiva est 1, 
\edtext{celeritas respectiva\indexS{celeritas respectiva}}{%
\lemma{celeritas}%
\Bfootnote{%
\textit{(1)}~cujusque %
\textit{(2)}~respectiva~\textit{L}}}
distributa in unoquoque $\displaystyle\frac{1}{2}$. Vis eorum respectiva\indexS{vis respectiva} est 
\rule[-2.4ex]{0pt}{6ex}%
$\displaystyle\frac{1}{4}C+\displaystyle\frac{1}{4}D=\displaystyle\frac{1}{2}$ 
quae addita priori vi respectivae\indexS{vis respectiva} $\displaystyle\frac{1}{4}$ facit 
\rule[-2.6ex]{0pt}{6ex}%
$\displaystyle\frac{3}{4}$
vim respectivam\indexS{vis respectiva} collectam, ut ante. Res ergo succedit. \pend 
\pstart Hinc sequitur posse motum totum corporis majoris transferri in minus.\indexS{translatio omnis motus} Ponamus %
\edtext{enim \textit{A} et \textit{B} globos}{\lemma{enim}\Bfootnote{%
\textit{(1)}~\textit{B} et %
\textit{(2)}~\textit{A} et \textit{B} %
\textit{(a)}~seca %
\textit{(b)}~esse %
\textit{(c)}~in \textit{\ind{3}} %
\textit{(d)}~globos~\textit{L}}} 
quiescere in locis \textit{\ind{3}A\ind{3}B}, et ex locis \textit{\ind{3}C},%
\edtext{\textit{\ind{3}D} venientes globos}{\lemma{\textit{\ind{3}D}}%
\Bfootnote{%
\textit{(1)}~venientia %
\textit{(2)}~venientes globos~\textit{L}}} 
\textit{D} et \textit{C} incurrere ipsis in \textit{\ind{2}C,\ind{2}D}, esse autem omnes aequales inter %
\edtext{se, globi \textit{C} et \textit{D},}{\lemma{se,}\Bfootnote{%
\textit{(1)}~quiescent tres globi \textit{B}, et \textit{C} et \textit{D} %
\textit{(2)}~globi \textit{C} et \textit{D},~\textit{L}}}
redigentur ad quietem\lbrack,\rbrack\ globus \textit{B} manebit in quiete, et globus \textit{A} recipiet omnem motum. Nam necesse %
est redire statum priorem. Nam eaedem %
\edtext{aequalita\lbrack te\rbrack s}{\lemma{}\Bfootnote{aequalitas \textit{L ändert Hrsg.}}} 
quae definiunt posterius ex priore, definiunt et contra, quia %
\edtext{aequationes hic semper sunt}{\lemma{aequationes}\Bfootnote{%
\textit{(1)}~sunt %
\textit{(2)}~hic semper sunt~\textit{L}}} 
rationales,\indexS{aequatio rationalis} seu non ambiguae,\indexS{aequatio ambigua} neque adeo dantur plures casus.
\pend 
\pstart Generaliter dicendum\lbrack:\rbrack\ si globi aequales\indexS{globi aequales}, quotcunque globos \textit{D} et \textit{C}, tangat globus \textit{B}, dum percutiatur ab \textit{A}, solo %
\edtext{moto, non progredi \textit{A}, quia non progreditur si solum \textit{B} percutiat,}{\lemma{moto,}\Bfootnote{%
\textit{(1)}~quiescere %
\textit{(2)}~quiescere %
\textit{(a)}~\textit{B}, quia %
\textit{(b)}~\textit{A}, %
\textit{(3)}~non progredi \textit{A}, %
\textit{(a)}~quia quiescit si solum \textit{B} percutiat, %
\textit{(b)}~quia non progreditur si solum %
\textit{(aa)}~\textit{A} 
\textit{(bb)}~\textit{B} percutiat,~\textit{L}}} 
ergo multo minus si \textit{B} cum %
\edlabel{35_13_03_242_4a}%
\edtext{}{%
{\xxref%
{35_13_03_242_4a}{35_13_03_242_4b}}%
\lemma{aliis.}%
\Bfootnote{%
\textit{(1)}~Porro \textit{B} 
\textit{(a)}~interpositus %
\textit{(b)}~globus inter %
\textit{(aa)}~tres %
\textit{(bb)}~plures positus aequales, comprimitur ab omnibus. Fingamus %
\textit{(aaa)}~esse %
\textit{(bbb)}~meras pilas inflatas moveri reciproce %
\textit{(2)}~Porro~\textit{L}}}%
aliis.
\pend
\begin{edfigure}%
\includegraphics[scale=1]{%
gesamttex/edit_VIII-4/images/LH_35_13_03_242_d3_242r.pdf} 
\edcaption{\edtext{\lbrack\textit{Fig.~3}\rbrack}{%
\lemma{\lbrack\textit{Fig.~3}\rbrack}%
\Cfootnote[nonum]{Ein gestrichener Entwurf zum Diagramm wird nicht wiedergegeben.}}}
\end{edfigure}
\pstart
Porro%
\edlabel{35_13_03_242_4b}
\textit{B}, sive solus %
\edtext{concurrat cum}{\lemma{concurrat}\Bfootnote{%
\textit{(1)}~, sive %
\textit{(2)}~cum~\textit{L}}}
globo \textit{A}, sive 
ab aliis prematur, vel %
\edtext{retineatur, videndum}{\lemma{retineatur,}\Bfootnote{%
\textit{(1)}~idem debet esse ictus %
\textit{(2)}~videndum~\textit{L}}}
in \textit{A}  
an eadem sit actio;\indexS{actio} id videtur dicendum cum \textit{A} idem %
\edtext{patiatur. Nempe concurrentes globi}{\lemma{patiatur.}\Bfootnote{%
\textit{(1)}~Itaque accuratius %
\textit{(2)}~Nota, quoad ictum %
\textit{(3)}~\bars{Considerandum \textit{streicht Hrsg.}} %
\textit{(4)}~Nempe %
\textit{(a)}~concurrens cum \textit{A}, e %
\textit{(b)}~concurrentes globi~\textit{L}}}
\textit{A} et \textit{B}, 
eodem modo percutiuntur aguntque et reagunt, sive adsint \textit{C}, \textit{D}, \textit{E}, sive non; itaque \textit{A} quiescit\lbrack,\rbrack\ \textit{B} recipit celeritatem ipsius \textit{A}, sed hac celeritate \textit{B} incurrens in \textit{C}, \textit{D}, \textit{E} aequales deberet quiescere, si haec regula succederet, %
\edtext{quod puto non esse verum.}{\lemma{}\Afootnote{\textit{Am Rand}: NB.}} 
Quaeritur generaliter, an res eodem modo fiant, sive ponas \textit{B} tangere ipsos \textit{C}, \textit{D}, \textit{E}, sive intervallum\indexS{intervallum quantumvis exiguum} %
\edtext{relinqui quantumvis}{\lemma{relinqui}\Bfootnote{%
\textit{(1)}~exiguum. Haec %
\textit{(2)}~quantumvis~\textit{L}}} 
exiguum. Haec regula si esset vera, haberetur quaesitum pro demonstrato. Et certe videtur admittenda: quia et circa %
\edtext{potentiarum\indexS{potentia} aequationes}{\lemma{potentiarum}\Bfootnote{%
\textit{(1)}~tam absolutarum quam re %
\textit{(2)}~aequationes~\textit{L}}} 
eadem utcunque exiguo discrimine\indexS{discrimen utcunque exiguum} provenire debent; scilicet quoad %
\edtext{vi\lbrack res\rbrack}{\lemma{}\Bfootnote{vis \textit{L ändert Hrsg.}}} 
respectivas,\indexS{vis respectiva} et progressivas,\indexS{vis progressiva} quomodocunque sumtas. 
\pend
\begin{edfigure}%
\includegraphics[scale=1]{%
gesamttex/edit_VIII-4/images/LH_35_13_03_242_d4_242r.pdf} 
\edcaption[0.25em]{\lbrack\textit{Fig.~4}\rbrack}
\end{edfigure}
\pstart
Experiamur\lbrack:\rbrack\ %
incurrat \textit{B} in \textit{C} et \textit{D}, et ponamus  
\textit{C} et \textit{D} moveri, ut ante. 
Ante ictum celeritas 1. 
Vis\indexS{vis} 1. 
Celeritas respectiva\indexS{celeritas respectiva} 1. 
Distributa\indexS{celeritas respectiva distributa} pro \textit{B}, 
\rule[-1.8ex]{0pt}{5.4ex}%
$\displaystyle\frac{2}{3}$; pro \textit{M}, 
\rule[-1.8ex]{0pt}{5.4ex}%
$\displaystyle\frac{1}{3}$. 
Ergo vis respectiva\indexS{vis respectiva} $4B+1M\lrcorner:9=%
\edtext{6:9=\displaystyle\frac{2}{3}$. Celeritas}{\lemma{$6:9=\displaystyle\frac{2}{3}$.}\Bfootnote{%
\textit{(1)}~Vis communis %
\textit{(2)}~Celeritas~\textit{L}}} 
centri gr.\indexS{celeritas centri gravitatis}\ est $1B+0M\lrcorner:\llcorner B+M=\displaystyle\frac{1}{3}$.%
\rule[-1.8ex]{0pt}{5.4ex} 
Ergo vis communis\indexS{vis communis} est 
\rule[-1.8ex]{0pt}{5.4ex}%
$\displaystyle\frac{3}{9}=\displaystyle\frac{1}{3}$, ut oportuit. 
Haec ante ictum,\indexS{ictus} post ictum quaeramus vim %
\edtext{communem.\indexS{vis communis} Celeritas centri gr.\indexS{celeritas centri gravitatis} $=}{\lemma{communem.}\Bfootnote{%
\textit{(1)}~Ea est %
\textit{(2)}~Celeritas centri gr.\ =~\textit{L}}} %
0B+\displaystyle\frac{1}{2}M\lrcorner:\llcorner B+M=\displaystyle\frac{1}{3}$ %
\edtext{ut ante.}{\lemma{}\Bfootnote{ut ante \textit{erg.~L}}} 
Ergo vis communis\indexS{vis communis} post 
\edtext{ictum est \rule[-1.8ex]{0pt}{5.4ex}%
$\displaystyle\frac{3}{9}=\displaystyle\frac{1}{3}$. Vis}{\lemma{ictum est}\Bfootnote{%
\textit{(1)}~$\displaystyle\frac{3}{36}=\displaystyle\frac{1}{\langle12\rangle}$. 
\bars{Sed post minor quam ante. \textit{erg.}} 
\textit{(2)}~$\displaystyle\frac{3}{9}=\displaystyle\frac{1}{3}$. Vis~\textit{L}}} 
respectiva\indexS{vis respectiva} inter \textit{C} et \textit{D} est 
\rule[-1.8ex]{0pt}{5.4ex}%
$\displaystyle\frac{1}{2}$ ut ante. 
Sed inter \textit{B} et \textit{M} celeritas respectiva\indexS{celeritas respectiva} 
est $\displaystyle\frac{1}{2}$, quae distributa,\indexS{celeritas respectiva distributa} pro \textit{B}, est 
\rule[-2.4ex]{0pt}{6ex}%
$\displaystyle\frac{2}{3}\;\displaystyle\frac{1}{2}=\displaystyle\frac{1}{3}$; pro \textit{M} est 
\rule[-2.4ex]{0pt}{6ex}%
$\displaystyle\frac{1}{3}\;\displaystyle\frac{1}{2}=\displaystyle\frac{1}{6}$. 
Fit vis respectiva\indexS{vis respectiva} 
\rule[-2.4ex]{0pt}{6ex}%
$\displaystyle\frac{4B+M}{36}=\displaystyle\frac{1}{6}$ cui adde $\displaystyle\frac{1}{2}$ fit 
\rule[-2.4ex]{0pt}{6ex}%
$\displaystyle\frac{1}{2}+\displaystyle\frac{1}{6}=\left.\displaystyle\frac{4}{6}\,\middle|\,\displaystyle\frac{2}{3}\right.$. Quae est tota vis respectiva %
\edtext{collecta.\indexS{vis respectiva collecta} Huic addita vis communis\indexS{vis communis}
$\displaystyle\frac{1}{3}$\rule[-1.8ex]{0pt}{5.4ex} reddit totam vim 1. Itaque procedit regula. Et habemus}{\lemma{collecta.}%
\Bfootnote{%
\textit{(1)}~Sed huic addita vis communis $\displaystyle\frac{1}{12}$ non producit totam vim 1. Itaque non procedit regula. Inquirendum esset in causas. Et video tantum interesse, quod hoc modo %
\textit{(2)}~Huic \lbrack...\rbrack\ habemus~\textit{L}}}
demonstrationem generalem %
\edtext{quaesiti. Et}{\lemma{quaesiti.}\Bfootnote{%
\textit{(1)}~Et regulam %
\textit{(2)}~Et~\textit{L}}} 
artificium prodit usus pulcherrimi ad multa difficilia solvenda. 
\pend
\pstart 
\edlabel{35_13_03_242_7a}%
Possumus hoc pronuntiare %
\edtext{theorema speciale:}{\lemma{theorema}\Bfootnote{%
\textit{(1)}~generale %
\textit{(2)}~speciale:~\textit{L}}}
si in tribus
\edtext{se\lbrack mi\rbrack quadrati}{%
\lemma{}%
\Bfootnote{%
sequiquadrati %
\textit{L ändert Hrsg.}}}
angulis sint centra trium globorum aequalium\indexS{globi tres aequales in tribus angulis semiquadrati} \textit{B}, \textit{C}, \textit{D}, et per latera quadrati %
\edtext{ven\lbrack i\rbrack entes}{\lemma{}\Bfootnote{venentes \textit{L ändert Hrsg.}}}
duo globi aequali celeritate percutiant tertium\indexS{globi duo in unum incurrentes} positum in laterum concursu, quiescent percutientes, et 
\edtext{percu\lbrack s\rbrack sus}{%
\lemma{}%
\Bfootnote{%
percursus %
\textit{L ändert Hrsg.}}}
perget in diagonali, celeritate quae 
\edtext{sit ad}{%
\lemma{sit}%
\Bfootnote{%
\textit{(1)}~ut %
\textit{(2)}~ad~\textit{L}}}
celeritatem incurrentium ut diagonalis ad latus. %
\edtext{Hoc pro Papino,\indexN{\textls{Papin} (Papinus), Denis 1647\textendash1713} %
ad ostendendum, %
\edtext{totam vim\indexS{vis}}{\lemma{totam}\Bfootnote{%
\textit{(1)}~vim globi %
\textit{(2)}~vim~\textit{L}}}
majoris materiae posse transferri\indexS{translatio totius vis majoris materiae in minorem} in 
\edlabel{35_13_03_242_3a}%
\edtext{}{{\xxref{35_13_03_242_3a}{35_13_03_242_3b}}%
\lemma{minorem.}\Bfootnote{%
\textit{(1)}~Unde %
\textit{(2)}~Videtur~\textit{L}}}%
minorem.\edlabel{35_13_03_242_7b}%
}{\lemma{Hoc \lbrack...\rbrack\  minorem}%
\Cfootnote{%
Siehe \indexN{\textls{Papin} (Papinus), Denis 1647\textendash1713}\textsc{D.~Papin},
\cite{00537}\textit{Mechanicorum de viribus motricibus sententia}, \cite{01023}\textit{AE} Januar 1691, S.\;6\textendash13, bes.\ S.\;9\textendash11.
Die von \indexN{\textls{Papin} (Papinus), Denis 1647\textendash1713}Papin im genannten Stoßfall angezweifelte These hatte Leibniz in \cite{02040}\textit{De causa gravitatis}, \cite{01023}\textit{AE} Mai 1690, S.\;228\textendash239 aufgestellt.}}
\pend
\pstart
Videtur%
\edlabel{35_13_03_242_3b}
corrigendus
\edtext{calculus meus prior quo putebam semper pergere incurrens.}{\lemma{calculus \lbrack...\rbrack\  incurrens}%
\Cfootnote{%
Anspielung noch nicht aufgelöst.}}
\pend
\pstart
\noindent
\edlabel{35_13_03_242_6a}%
\href{\Scan{LH_35_13_03_242v}}%
{\lbrack242~v\textsuperscript{o}\rbrack}
\pend
\nobreak
\begin{edfigure}[0pt]%
\includegraphics[scale=1]{%
gesamttex/edit_VIII-4/images/LH_35_13_03_242_d5_242v.pdf}
\edcaption{\lbrack\textit{Fig.~5}\rbrack}
\end{edfigure}
\pstart 
Res bene succedit in casu illo speciali, sed nescio an generaliter pro incursu corporis in duo\indexS{incursus corporis in duo}. Ex principiis nostris, si vera est %
\edtext{Regula Hugeniana\indexS{regula Hugeniana} %
\edtext{ut globus \textit{A} incurrens in \textit{B}, tangentem plures\indexS{globus plures globos tangens}}{\lemma{ut}\Bfootnote{%
\textit{(1)}~corpus \textit{A} incurrens in \textit{B}, tangens plura %
\textit{(2)}~globus \lbrack...\rbrack\ plures~\textit{L}}} %
\textit{C}, \textit{D}, \textit{E} quiescat cum \textit{B},%
}{\lemma{Regula \lbrack...\rbrack\ cum \textit{B}}%
\Cfootnote{%
\indexN{\textls{Huygens} (Hugenius, Ugenius, Hugens, Huguens), Christiaan 1629\textendash1695}\textsc{C.~Huygens}, \cite{00525}\textit{Traité de la lumière}, Leiden 1690, chap.~1, S.\;14 (\textit{HO}\cite{00113} XIX, S.\;473). %
Siehe Leibnizens Auszüge aus dem einschlägigen Abschnitt des \cite{00525}\textit{Traité}:
\cite{02094}\textit{LSB} VIII,~3 N.\;75 von September\textendash Oktober 1690.}}
sequetur idem esse si \textit{B} incurrat in \textit{C}, \textit{D}, \textit{E}, quia perinde est ac si \textit{B} non tangat, sed parum a tangente absit, eo autem casu quiescit \textit{A}, translato motu in \textit{B} %
\edtext{quod incurrat in \textit{C}, \textit{D}, \textit{E},}{\lemma{}\Bfootnote{quod incurrat in \textit{C}, \textit{D}, \textit{E} \textit{erg.~L}}} 
ergo ut prodeant, quae ante, debet \textit{B} incurrens in \textit{C}, \textit{D}, \textit{E}, quietem accipere. Et %
\edtext{hoc de duobus}{\lemma{hoc}\Bfootnote{%
\textit{(1)}~duo %
\textit{(2)}~de duobus~\textit{L}}}
verum reperi, si sint %
\edtext{siti invicem}{\lemma{siti}\Bfootnote{%
\textit{(1)}~in se %
\textit{(2)}~invicem~\textit{L}}} 
in Triangulo rectangulo isoscele. Sed secus inquirendum de tribus certe, %
\edtext{an Hugenianae regulae\indexS{regula Hugeniana} obstat}{\lemma{an}\Bfootnote{%
\textit{(1)}~vera sit %
\textit{(2)}~Hugenianae regulae %
\textit{(a)}~obstant %
\textit{(b)}~obstat~\textit{L}}} 
\edlabel{35_13_03_242_5a}%
\edtext{}{%
{\xxref%
{35_13_03_242_5a}{35_13_03_242_5b}}%
\lemma{ratiocinatio \lbrack...\rbrack\ desinentiarum}%
\Cfootnote{%
Zum Kontinuitätsprinzip (\glqq regula desinentiarum\grqq, auch \glqq regula continuitatis\grqq) vgl.\ 
\textsc{G.\,W.~Leibniz}, 
\cite{02054}\textit{Extrait d'une Lettre de M.~L.\ sur un Principe Général}, in 
\cite{02002}\textit{NRL} Juli 1687, S.\;744\textendash753 
(\cite{02039}\textit{LPG} III, S.\;51\textendash55); sowie 
\cite{02055}\textit{Principium quoddam generale}, 
\textit{LSB} VI, 4, N.\;371.}}%
ratiocinatio firma meae
\edtext{regulae desinentiarum.\indexS{regula desinentiarum}\edlabel{35_13_03_242_5b}}{\lemma{regulae}\Bfootnote{%
\textit{(1)}~de Desinentiis %
\textit{(2)}~seu reg %
\textit{(3)}~desinentiarum.~\textit{L}}}
Nempe,
\edtext{globus}{%
\lemma{}%
\Bfootnote{%
globus %
\textit{erg.~L}}}
\textit{A} incurrat simul in globos \textit{B} et \textit{D}. Et propemodum tangat globum \textit{C}, dico eventus qui fiant esse debere tales, ut si tandem status hic desinat in statum perfecti contactus\indexS{contactus perfectus} etiam ipsius \textit{C}, etiam eventus in hunc desinant. 
Quod verum non est, si vera regula Hugeniana\indexS{regula Hugeniana}, nam si tantillum absit, %
\edtext{quiescet \textit{A}, procedent \textit{B},}{\lemma{quiescet \textit{A},}\Bfootnote{%
\textit{(1)}~prodibunt \textit{B} %
\textit{(2)}~procedent \textit{B},~\textit{L}}} %
\edtext{et \textit{D}, manente}{\lemma{et \textit{D},}\Bfootnote{%
\textit{(1)}~quiescente %
\textit{(2)}~manente~\textit{L}}}
in quiete ipso \textit{C}, ergo non desinitur %
\edtext{in regulae}{\lemma{in}\Bfootnote{%
\textit{(1)}~statum %
\textit{(2)}~regulae~\textit{L}}} 
Hugenianae\indexS{regula Hugeniana} eventum, %
\edtext{quem esse}{\lemma{quem}\Bfootnote{%
\textit{(1)}~vult %
\textit{(2)}~esse~\textit{L}}} 
si \textit{A} tres \textit{B}, \textit{C}, \textit{D} simul tangat sequitur ex ejus regula, ut scil.\ tunc quiescat \textit{A}, procurrant \textit{B}, \textit{C}, \textit{D}.
\pend
\begin{edfigure}%
\includegraphics[scale=1]{%
gesamttex/edit_VIII-4/images/LH_35_13_03_242_d6_242v.pdf} 
\edcaption[0.25em]{\edtext{\lbrack\textit{Fig.~6}\rbrack}{%
\lemma{\lbrack\textit{Fig.~6}\rbrack}%
\Cfootnote[nonum]{Ein gestrichener Entwurf zum Diagramm wird nicht wiedergegeben.}}}
\end{edfigure}
\pstart 
Imo contra regulam Hugenianam\indexS{regula Hugeniana} per %
\edtext{se sumtam (\protect\vphantom)sine}{\lemma{se sumtam}\Bfootnote{%
\textit{(1)}~aliter %
\textit{(2)}~(\protect\vphantom)sine~\textit{L}}}
consequentia\protect\vphantom() sic pugnat 
regula desinentiarum\indexS{regula desinentiarum} seu Regula continuitatis\indexS{regula continuitatis}.
Nam si semper verum manet, ut \textit{A} incurrente in \textit{B} tangentem \textit{C} et \textit{D}, quiescant \textit{A} et \textit{B}, translato motu in \textit{C} et \textit{D}, sequitur et hoc futurum si angulus \textit{CBD} sit  %
\edtext{infinite \lbrack obtusus\rbrack,\indexS{angulus infinite obtusus} quo}{\lemma{infinite}\Bfootnote{\bars{acutus \textit{ändert Hrsg.}} %
\textit{(1)}~quo 
\textit{(2)}~\bars{et \textit{streicht Hrsg.}} %
\textit{C}, \textit{D} tantillum 
\textit{(3)}~quo~\textit{L}}} 
casu falsum est \textit{B} quiescere omni motu in \textit{C} et \textit{D} translato. 
Quod si quis vellet hoc casu quo scil.\ \textit{CBD} sunt in eadem recta tunc \textit{B} %
\edtext{ab \textit{A} impulsum}{\lemma{ab \textit{A}}\Bfootnote{%
\textit{(1)}~in qua %
\textit{(2)}~impulsum~\textit{L}}}
quiescere, %
\edtext{sequetur non}{\lemma{sequetur}\Bfootnote{%
\textit{(1)}~\textit{C} et \textit{D} %
\textit{(2)}~non~\textit{L}}} 
progredi centrum gravitatis totale\indexS{centrum gravitatis totale}, nam ipsis \textit{A} et \textit{B} quiescentibus, et ipsis \textit{C} et \textit{D} motum recipientibus et quidem in recta ipsa \textit{CD}, %
\edtext{patet nihil}{\lemma{patet}\Bfootnote{%
\textit{(1)}~non %
\textit{(2)}~nihil~\textit{L}}}
\edtext{u\lbrack lt\rbrack ra}{\lemma{}\Bfootnote{utlra \textit{L ändert Hrsg.}}} 
hanc rectam prodire, atque adeo centrum gravitatis\indexS{centrum gravitatis} sisti.%
\edlabel{35_13_03_242_6b}%
\pend
